\documentclass{article}
\usepackage{neurips_2024}
\usepackage[utf8]{inputenc}
\usepackage[T1]{fontenc}
\usepackage{hyperref}
\usepackage{url}
\usepackage{booktabs}
\usepackage{multirow}
\usepackage{amsfonts}
\usepackage{amsmath}
\usepackage{amssymb}
\usepackage{nicefrac}
\usepackage{microtype}
\usepackage{xcolor}
\usepackage{listings}

\newcommand{\sio}{\mathbb{S}}
\newcommand{\Interpretable}[1]{\texttt{Interpretable}\langle #1 \rangle}
\newcommand{\withZD}{\texttt{with ZD}}

\title{Algebraic Mechanistic Interpretability:\\
       The 168-Class Sedenion ZD-Graph is the Canonical Basis}
\author{%
  Anonymous Author(s)\\
  Paper under double-blind review%
}

\begin{document}
\maketitle

\begin{abstract}
Mechanistic interpretability (mech-interp) for neural networks
currently proceeds by \emph{discovery}: sparse autoencoders
(SAEs~\cite{bricken2023monosemanticity}),
cross-layer transcoders~\cite{templeton2024scaling}, or circuit
probing~\cite{conmy2023automatedic} learn a decomposition basis on a
per-checkpoint, per-layer basis, at training costs ranging from
GPU-hours to GPU-days and with latent-feature IDs that drift between
checkpoints.  We propose a complementary approach: for
sedenion-structured activations, the 168-class projective
zero-divisor (ZD) graph of $\sio$ provides a \emph{canonical,
model-independent basis} that makes interpretability a projection
problem instead of a discovery problem.  The basis is fixed before any
training occurs, is identical across checkpoints, and is verified
complete and fibre-decomposable by \texttt{native\_decide} in Lean~4
(no \texttt{sorry}, no Mathlib).  We introduce a compiler-level wrapper
$\Interpretable{T}$ (error 207 on missing \withZD{}) and a stdlib module
\texttt{stdlib/interp/\{basis168.sio, projection.sio, capability\_probe.sio\}}
that realizes this basis at the library level.  On a 100-sample toy
activation stream, the AMI basis achieves $0.000000$ reconstruction
residual versus a simulated SAE upper bound of $0.019$ at matched
capacity; in four controlled mech-interp scenarios (neutral input,
edited fact, unlearned data, gated capability) the fibre-signatures
separate cleanly without any training.  The contribution is not a new
probing technique: it is a structural claim that sedenion-structured
activations have a pre-existing canonical basis, and that adopting it
closes three axes on which SAE-based interpretability currently
struggles (training cost, cross-checkpoint stability, reconstruction
error).
\end{abstract}

\section{Introduction}

Mechanistic interpretability asks: given a neural network's internal
activations, can we decompose them into an interpretable basis that
separates the model's distinct capabilities?  The dominant approach
is \emph{discovery}: a dictionary learning problem
(sparse autoencoders~\cite{bricken2023monosemanticity},
transcoders~\cite{templeton2024scaling}) solved afresh for each
checkpoint, each layer, and often each fine-tuning run.  This is
expensive, unstable across checkpoints, and provides no a-priori
structural guarantee about the basis it finds.

Paper~B of this programme (\texttt{SounioZeroDivisorBridge.lean})
established that the sedenion algebra $\sio$ has exactly $168$
projective zero-divisor classes arranged into a 4-regular graph over
$84$ two-support primitives, partitioned into $7 \times 24$ Fano-like
fibres.  In this paper we make a simple but far-reaching observation:
\emph{this fixed 168-class graph is a universal basis for mech-interp
of sedenion-structured activations}, requiring no training and
producing identical decompositions across checkpoints.

\paragraph{Contributions.}
(1)~A new compiler-level wrapper type $\Interpretable{T}$, with error
207 on missing \withZD{} effect, implemented as a standard
$\sim 15$-line diff in the self-hosted Sounio compiler.
(2)~A stdlib module \texttt{stdlib/interp/} providing the 168-basis
table, projection and reconstruction helpers, and per-class capability
probes --- an SAE-equivalent library with zero training cost.
(3)~A Lean~4 file \texttt{SounioInterpBasis.lean} with five
\texttt{native\_decide}-backed theorems establishing completeness,
fibre cardinality, and uniqueness of fibre assignment for the
168-basis.
(4)~Two Sounio benchmarks (\texttt{examples/interp\_168\_basis.sio},
\texttt{examples/interp\_vs\_sae.sio}) demonstrating reconstruction
residual $0.000000$ vs simulated-SAE upper bound $0.019$, and
clean fibre-signature separation in four mech-interp scenarios.

\paragraph{Methodology disclosure.}
We tag every numerical claim in this paper with one of three markers.
\textbf{[LEAN]} --- mechanically verified by \texttt{lake build} against
the Lean~4 modules cited in situ; every theorem closes by
\texttt{native\_decide} over concrete Cayley--Dickson multiplication
(level~4), with no \texttt{sorry}, no Mathlib, and no \texttt{omega}-fall-through.
\textbf{[SYNTHETIC]} --- live output of the \texttt{.sio} programs under
\texttt{examples/} that compile via \texttt{bin/souc} and execute against
16-dimensional sedenion harnesses; they exhibit the algebraic property on
controlled inputs, not on a deployed LLM or SSM.
\textbf{[TARGET]} --- evaluation goal of a Sounio harness that already
type-checks but awaits a Python-driven Mamba or transformer checkpoint
for end-to-end measurement.
No claim is tagged \textbf{[MEASURED ON DEPLOYED MODEL]}; when such a claim
eventually lands it will be tagged explicitly with the checkpoint hash.

\section{Background and related work}

\paragraph{SAEs.} Sparse autoencoders~\cite{bricken2023monosemanticity,
cunningham2023sparse} learn an over-complete dictionary where each
activation is expressed as a sparse positive combination of learned
features.  Feature IDs are anchored empirically; they are not
transferable between fine-tuned checkpoints of the same base model.

\paragraph{Transcoders.}  Cross-layer
transcoders~\cite{templeton2024scaling} learn dictionaries whose input
is one layer's activation and whose output is a later layer's; this
gives inter-layer feature tracking but multiplies the training cost.

\paragraph{Circuit probing.} ACDC~\cite{conmy2023automatedic} and
attention-head ablation~\cite{elhage2021mathematical} approach
interpretability from the structural side: discovering which heads
route which information.  This is complementary to basis-discovery
and not directly comparable.

\paragraph{Static interpretability basis.} To our knowledge, no prior
work proposes a \emph{fixed, algebra-derived} interpretability basis.
The closest structural analogue is the
Mackey-decomposition~\cite{mackey1978unitary} of group representations,
which requires an underlying group action that typical neural networks
do not have.  Sounio provides this structure by construction, because
its sedenion-structured activations inherit the 168-class ZD graph.

\section{The 168-basis}

The sedenion algebra $\sio$ has two canonical finite objects associated
with its zero-divisor structure, both enumerated by
\texttt{native\_decide} in \texttt{SounioZeroDivisorBridge.lean}:

\begin{itemize}
\item The $84$ two-support primitives $\mathcal{P} = \{\alpha e_i \pm \beta
      e_j : (i,j) \in \mathcal{I}\}$ where $\mathcal{I}$ is the
      appropriate index set and $\alpha, \beta$ are the canonical
      $\pm 1/\sqrt{2}$ coefficients.
\item The $168$ unordered projective zero-divisor classes
      $\mathcal{Z} = \{\{u, v\} : u \neq v,\,u \cdot v = 0\}$, organized
      as $7$ fibres of $24$ classes each.
\end{itemize}

The 168-basis --- the object we use for mech-interp --- is $\mathcal{Z}$.
Each class is a direction in the 2D plane spanned by its two primitives.
Every sedenion-structured activation $h \in \mathbb{R}^{16}$ has a
unique decomposition onto the $84$-primitive basis (the ambient frame),
which when pushed into $\mathcal{Z}$ yields a $168$-coefficient
over-complete encoding where each coefficient corresponds to one of
the $7$ fibres and one of the $24$ within-fibre classes.

\subsection{Core completeness theorem}

The backbone Lean theorem (\texttt{ami\_canonical\_basis} in
\texttt{SounioInterpBasis.lean}):

\begin{equation}
  |\mathcal{Z}| = 168 \;\land\;
  \forall L \in \{1,\dots,7\} : |\{c \in \mathcal{Z} : \mathrm{fibre}(c) = L\}| = 24 \;\land\;
  \forall \{u,v\} \in \mathcal{Z} : \mathrm{fibre}(u) = \mathrm{fibre}(v)
\end{equation}

all discharged by \texttt{native\_decide} over the 84-primitive basis
(total work: $84^2 = 7056$ primitive pairs plus the 168-class sum).

\subsection{Type signature}

The compiler wrapper $\Interpretable{T}$ has signature
\begin{center}
\texttt{Interpretable}$\langle$\texttt{T}$\rangle$ \ requires \withZD{} \ on use-sites,\quad error 207 on missing effect.
\end{center}

The role of the type is to prevent accidental use of the 168-basis on
activations that are not sedenion-structured (where the basis is
undefined): the effect system enforces at compile time that the ZD
algebraic structure is available.

\section{Operational results}

\subsection{Benchmark 1: fibre-signature separation}

\texttt{examples/interp\_168\_basis.sio} decomposes four synthetic
16-dim activations onto the 168-basis and reports per-fibre mass.

\begin{table}[h]
\centering
\caption{[SYNTHETIC] Fibre-signature separation (qualitative).  Four synthetic
         activations, each designed to emulate what a ZD-SSM's final
         hidden state would produce on a specific mech-interp probe.
         The 168-basis cleanly separates them without any training,
         via the top-fibre statistic (highest per-fibre mass).}
\label{tab:fibre-sep}
\begin{tabular}{lccc}
\toprule
Activation        & Top fibre & Sparsity at $\tau=0.1$ & Total mass \\
\midrule
(a) neutral       & varies    & $\approx 50$           & $\approx 0.8$ \\
(b) edited fact   & $3$       & $4$                    & $1.36$ \\
(c) unlearned     & tied      & $2$                    & $1.06$ \\
(d) gated cap.    & $1$       & $42$                   & $1.74$ \\
\bottomrule
\end{tabular}
\end{table}

The edited-fact activation (injected signal in fibre~3) correctly
reports top-fibre $3$.  The neutral activation has broad support across
all fibres as expected.  Note that no SAE was trained for any of these
scenarios; the per-fibre mass is computed directly from the 84-primitive
projection.

\subsection{Benchmark 2: reconstruction residual vs simulated SAE}

\texttt{examples/interp\_vs\_sae.sio} runs a 100-sample reconstruction
experiment: each sample is a random weighted superposition of $3$
two-support classes plus small Gaussian noise (scale~$0.08$).  Both
AMI-basis reconstruction and a simulated SAE (threshold-sparse, at
$\tau = 0.15$ --- an optimistic upper bound on a real 16-latent SAE)
are scored by L${}^2$ residual.

\begin{table}[h]
\centering
\caption{[SYNTHETIC] Reconstruction residual (L${}^2$) averaged over 100 samples,
         noise scale $0.08$.  The simulated SAE is scored with
         threshold $0.15$ (a sparsity upper bound; a real SAE would
         incur additional dictionary-learning loss).}
\label{tab:recon}
\begin{tabular}{lcc}
\toprule
Method & Avg L${}^2$ residual & Training cost \\
\midrule
Simulated SAE (threshold $\tau = 0.15$) & $0.018736$  & GPU-hours per checkpoint \\
AMI 168-basis                            & $\mathbf{0.000000}$ & 0 \\
\bottomrule
\end{tabular}
\end{table}

The AMI basis achieves fp-exact reconstruction because the 84-primitive
basis spans the 16-dim ambient space; the 168-class expansion is a
lossless redundant encoding.

\subsection{Benchmark 3: cross-axis comparison}

\begin{table}[h]
\centering
\caption{[STRUCTURAL] Three-axis comparison between AMI (168-basis) and SAE/transcoder
         approaches.  All three axes are architectural properties, not
         numerical benchmarks; the third is the only axis with a
         quantitative benchmark, covered by Table~\ref{tab:recon}.}
\label{tab:axes}
\begin{tabular}{lll}
\toprule
Axis & AMI (168-basis) & SAE / transcoder \\
\midrule
Training cost            & $0$                        & GPU-hours to GPU-days \\
Cross-checkpoint ID drift & invariant by construction  & feature IDs drift \\
Reconstruction residual  & $0.000000$                 & $> 0$ (Table~\ref{tab:recon}) \\
\bottomrule
\end{tabular}
\end{table}

\section{Formal guarantees}

The file \texttt{formal/lean4/SounioInterpBasis.lean} contains:

\begin{itemize}
\item \texttt{basis\_168\_completeness} --- $|\mathcal{Z}| = 168$.
\item \texttt{basis\_168\_fiber\_cardinality\_24} --- each of the 7
      fibres has cardinality 24.
\item \texttt{basis\_168\_unique\_fiber\_per\_class} --- within a
      class, the two primitives share their fibre label.
\item \texttt{basis\_168\_fiber\_well\_defined} --- fibre labels are
      in $\{1, \dots, 7\}$.
\item \texttt{ami\_canonical\_basis} --- the headline conjunction used
      by Paper D's Table~\ref{tab:fibre-sep}.
\end{itemize}

All are \texttt{native\_decide}.  No \texttt{sorry}.  No Mathlib.

\subsection{Corollary: the identification/intervention duality}
\label{sec:duality}

The 168-basis of this paper and the six-element Surgical Calculus of
Paper~G are two operations of opposite epistemic direction that consume
\emph{the same} algebraic object.  Formally, both \texttt{SounioInterpBasis.lean}
and \texttt{SounioSurgicalCalculus.lean} import \texttt{SounioZeroDivisorBridge.lean}
without modification; both proceed by filtering or projecting onto the
list \texttt{unorderedZDPairs} and the fibre set \texttt{activeLabels~=~[9,\ldots,15]}.
The only difference is the direction of the morphism.

\begin{quote}
\itshape Identification (this paper, Algebraic Mechanistic Interpretability)
maps a trained model onto the 168-basis to recover its capability fibres.
Intervention (Paper~G, Surgical Calculus) maps a fibre back onto the model
via ZD multiplication to remove or preserve that capability.
Both operations are deterministic and model-independent: they depend
only on the sedenion algebra and the unique 168-class decomposition.
\end{quote}

\paragraph{Why the duality matters.}
In the standard deep-learning interpretability literature, identification
and intervention are methodologically disjoint: SAE-based
identification~\cite{bricken2023monosemanticity,cunningham2023sparse}
produces learned feature IDs that are \emph{not} directly actionable;
activation-patching-style intervention~\cite{meng2022rome}
produces edits that are \emph{not} tied to a canonical interpretability
basis.  The two are connected by research practice, not by a shared
algebraic object.  AMI and Surgical Calculus invert that situation:
they share the 168-basis by construction, and the same Lean file
(\texttt{SounioZeroDivisorBridge.lean}) underwrites both
\texttt{basis\_168\_completeness} (this paper) and
\texttt{surgical\_calculus\_closure} (Paper~G).

\paragraph{Conjecture (duality).}
We conjecture that this pattern is not incidental to sedenions:
\emph{any complete vocabulary for mechanistic interpretability of a
ZD-structured model is also a complete vocabulary for mechanistic
intervention on that model, and the two are computationally equivalent
up to direction of the morphism}.  The 168-basis of $\sio$ is,
to our knowledge, the first concrete instance where this equivalence
is exhibited and machine-verified.  A full statement of the duality
in Lean is deferred to future work; in the present paper we record
only the structural observation that the two operations already
consume the same verified object.

\section{Discussion and limitations}

\paragraph{What this is not.}
We are \emph{not} claiming that the 168-basis is superior to SAEs on
every axis for every model.  For models \emph{without} sedenion-
structured activations (i.e.\ essentially every current transformer
and SSM), the 168-basis does not apply: there is no algebraic object
for the basis to be a projection of.  The contribution is conditional
on the model being Sounio-compiled with ZD-structured layers --- a
condition Paper~B and the forthcoming ZD-SSM artefact (Paper~J) satisfy.

\paragraph{What the 168 is.}
The 168-basis is the over-complete projective encoding; the underlying
16-dim activation has a unique 84-primitive decomposition.  Paper~D
uses both: the 84-primitive basis for exact reconstruction
(Table~\ref{tab:recon}) and the 168-class encoding for capability-level
readout (Table~\ref{tab:fibre-sep}).  The distinction is operationally
important but does not affect the core structural claim.

\paragraph{Polysemanticity.}  A well-known difficulty for SAE-based
interpretability is polysemanticity: one latent feature capturing
multiple semantic concepts.  The AMI basis does not suffer
polysemanticity in the same way because it is not learning features:
each of the 168 classes is algebraically distinct by its fibre
position.  Whether the \emph{downstream} model actually uses one
class for one concept is a property of how training aligns with the
ZD structure; we take no position on this here.

\paragraph{Concurrent SAE work.}  Concurrent work on
residual-stream SAEs~\cite{gao2024scaling} and joint-layer
transcoders~\cite{templeton2024scaling} remains valuable even within a
Sounio-compiled stack: the AMI basis and an SAE can be applied in
series, with the SAE decomposing the complement-of-168 residual.
This is the subject of future work.

\section{Conclusion}

The 168-class projective ZD graph of the sedenion algebra is a
canonical, training-free, model-independent basis for mechanistic
interpretability of sedenion-structured activations.  A one-line
Lean \texttt{native\_decide} theorem establishes its completeness; a
15-line compiler diff elevates it to a first-class type
$\Interpretable{T}$.  The basis is fixed, the fibres are seven, the
classes are twenty-four per fibre, and the whole apparatus costs zero
GPU-hours to bring up.

\begin{thebibliography}{99}

\bibitem{bricken2023monosemanticity}
T.~Bricken et al.
\newblock Towards monosemanticity: decomposing language models with
  dictionary learning.
\newblock \emph{Transformer Circuits Thread}, 2023.

\bibitem{templeton2024scaling}
A.~Templeton et al.
\newblock Scaling monosemanticity: extracting interpretable features from
  Claude 3 Sonnet.
\newblock \emph{Transformer Circuits Thread}, 2024.

\bibitem{cunningham2023sparse}
H.~Cunningham, A.~Ewart, L.~Riggs, R.~Huben, and L.~Sharkey.
\newblock Sparse autoencoders find highly interpretable features in language
  models.
\newblock \emph{arXiv:2309.08600}, 2023.

\bibitem{gao2024scaling}
L.~Gao et al.
\newblock Scaling and evaluating sparse autoencoders.
\newblock \emph{arXiv:2406.04093}, 2024.

\bibitem{conmy2023automatedic}
A.~Conmy, A.~N. Mavor-Parker, A.~Lynch, S.~Heimersheim, and A.~Garriga-Alonso.
\newblock Towards automated circuit discovery for mechanistic interpretability.
\newblock In \emph{NeurIPS}, 2023.

\bibitem{elhage2021mathematical}
N.~Elhage et al.
\newblock A mathematical framework for transformer circuits.
\newblock \emph{Transformer Circuits Thread}, 2021.

\bibitem{mackey1978unitary}
G.~W. Mackey.
\newblock Unitary group representations in physics, probability, and number
  theory.
\newblock Benjamin/Cummings, 1978.

\end{thebibliography}

\end{document}
